\begin{mistake}{2026-06-03}{04}

\begin{problemcontent}
设 \(\lim_{x \to 0} \varphi(x) = 0\)，则下列命题中正确的个数为
(1) \(\lim_{x \to 0} \frac{\sin \varphi(x)}{\varphi(x)} = 1\).
(2) \(\lim_{x \to 0} (1 + \varphi(x))^{\frac{1}{\varphi(x)}} = e\).
(3) 若 \(f'(x_0) = A\)，则 \(\lim_{x \to 0} \frac{f(x_0 + \varphi(x)) - f(x_0)}{\varphi(x)} = A\).
(4) 若 \(\lim_{u \to 0} f(u) = A\)，则 \(\lim_{x \to 0} f(\varphi(x)) = A\).

(A) 0 个. (B) 2 个. (C) 3 个. (D) 4 个.
\end{problemcontent}

\begin{solutioncontent}
正确答案为 (A) 0 个。

对于 (1)、(2)、(3)：题目未限制 \(\varphi(x) \neq 0\)。若取 \(\varphi(x) \equiv 0\)（满足 \(\lim_{x \to 0} \varphi(x) = 0\)），则三者的表达式分母均为 0，无意义，无法谈论极限。故 (1)、(2)、(3) 均错。

对于 (4)：考察复合函数极限。若取 \(\varphi(x) \equiv 0\)，且令外层函数 \(f(u)\) 满足 \(f(0) = B\) (\(B \neq A\))，且当 \(u \neq 0\) 时 \(f(u) = A\)。则 \(\lim_{u \to 0} f(u) = A\) 满足，但 \(\lim_{x \to 0} f(\varphi(x)) = \lim_{x \to 0} f(0) = B \neq A\)。故 (4) 错。

四个命题全错。
\end{solutioncontent}

\mistakeknowledge{
\begin{itemize}
  \item \textbf{核心陷阱}：变量替换时（特别是作分母时），必须保证在极限去心邻域内该变量不为 0。
  \item \textbf{复合函数极限定理}：\(\lim f(\varphi(x)) = \lim f(u)\) 的成立前提是外层函数 \(f\) 连续，或内层函数 \(\varphi(x) \neq u_0\)。
\end{itemize}}

\end{mistake}
